\documentclass[11pt,a4paper]{article}

\usepackage[utf8]{inputenc}
\usepackage[T1]{fontenc}
\usepackage[ngerman]{babel}
\usepackage{lmodern}
\usepackage{amsmath,amssymb,mathtools}
\usepackage{booktabs,array}
\usepackage{enumitem}
\usepackage{geometry}
\usepackage{xcolor}
\usepackage[most]{tcolorbox}
\usepackage{tikz}
\usepackage{hyperref}

\geometry{left=25mm,right=25mm,top=24mm,bottom=27mm}
\setlength{\parindent}{0pt}
\setlength{\parskip}{0.6em}
\setlist[itemize]{itemsep=0.2em,topsep=0.2em}
\setlist[enumerate]{itemsep=0.2em,topsep=0.2em}
\renewcommand{\arraystretch}{1.25}

\newcommand{\Bel}{\operatorname{Bel}}
\newcommand{\pBel}{\overline{\operatorname{Bel}}}
\newcommand{\LM}{\mathrm{L}}

\newtcolorbox{aufgabenstellung}{
  enhanced,
  breakable,
  colback=gray!13,
  colframe=gray!13,
  boxrule=0pt,
  arc=1mm,
  left=2mm,
  right=2mm,
  top=1mm,
  bottom=1mm
}

\title{Lösung zu Aufgabe B6.1 (Markow-Lokalisierung)}
\author{Richard Bethke, Leander Engel, Sven Scherzer, Katja Alt}
\date{02.07.2026}

\begin{document}
\maketitle

\begin{aufgabenstellung}
\textbf{Aus der Aufgabenstellung.}
Gegeben sei die abgebildete Welt eines mobilen Roboters. Der Roboter bewege sich deterministisch
in einem Korridor mit $16$ Gitterzellen. In einigen dieser Zellen sind Landmarken angebracht.
Wenn der Roboter sich in einer der Zellen mit Landmarken befindet, so wird diese Landmarke von dem
Roboter mit einer Wahrscheinlichkeit von $70\%$ detektiert. In allen anderen Zellen (ohne
Landmarken) meldet der Roboter aufgrund der Sensorungenauigkeit mit $25\%$ Wahrscheinlichkeit
dennoch eine Landmarke.

Der Roboter führt die folgenden $5$ Aktionen aus:
\begin{enumerate}
  \item Der Roboter detektiert eine Landmarke.
  \item Der Roboter bewegt sich $2$ Zellen im Uhrzeigersinn.
  \item Der Roboter detektiert wieder eine Landmarke.
  \item Der Roboter bewegt sich $4$ Zellen im Uhrzeigersinn.
  \item Der Roboter detektiert keine Landmarke.
\end{enumerate}
Berechnen Sie die Aufenthaltswahrscheinlichkeit für den Roboter für jede Zelle nach jedem
einzelnen Schritt. In welcher Zelle ist demnach der Roboter vermutlich gestartet?

Nehmen Sie nun an, dass bei einer Bewegung des Roboters zwar bekannt ist, wie weit sich der
Roboter bewegt hat, nicht aber in welche Richtung (im oder entgegen dem Uhrzeigersinn). Berechnen
Sie auch hier die Aufenthaltswahrscheinlichkeiten.
\end{aufgabenstellung}

Der ringförmige Korridor mit den Landmarkenzellen $\{1,4,6,9,13\}$:

\begin{center}
\begin{tikzpicture}[x=1cm,y=1cm]
  % Ringzellen -- die Mitte bleibt leer, sodass nur ein Kreis entsteht
  \foreach \c/\r/\lab in {%
    0/0/1, 1/0/2, 2/0/3, 3/0/4, 4/0/5,
    0/1/16, 4/1/6,
    0/2/15, 4/2/7,
    0/3/14, 4/3/8,
    0/4/13, 1/4/12, 2/4/11, 3/4/10, 4/4/9}{%
    \draw (\c,-\r) rectangle ++(1,-1);
    \node at (\c+0.5,-\r-0.5) {$\lab$};
  }
  % Landmarken als schwarze Quadrate in der Zellenecke
  \foreach \c/\r in {0/0, 3/0, 4/1, 0/4, 4/4}{%
    \fill (\c,-\r) rectangle ++(0.2,-0.2);
  }
\end{tikzpicture}
\end{center}
{\small (Die schwarzen Quadrate markieren die Zellen mit Landmarken $\{1,4,6,9,13\}$.)}

Gegeben sind die folgenden Wahrscheinlichkeiten:

\begin{center}
\setlength{\tabcolsep}{10pt}
\begin{tabular}{c|cc}
\toprule
 & $x=\LM$ & $x=\neg\LM$ \\
\midrule
$z=\LM$       & $0.7$ & $0.25$ \\
$z=\neg\LM$   & $0.3$ & $0.75$ \\
\bottomrule
\end{tabular}
\end{center}

\section*{Berechnung der Aufenthaltswahrscheinlichkeit (mit vorgegebener Richtung)}

Formeln:
\[
  \pBel(x_{t+1}) = \sum_{x_t} P(x_{t+1}\mid x_t,u_t)\cdot \Bel(x_t)
  \qquad\text{(a-priori-Schätzung)},
\]
\[
  \Bel(x_{t+1}) = \eta^{-1}\cdot P(z_t\mid x_{t+1})\cdot \pBel(x_{t+1})
  \qquad\text{(a-posteriori-Schätzung)},
\]
\[
  \qquad
  \eta^{-1} = \frac{1}{\sum_{x_{t+1}} P(z_t\mid x_{t+1})\cdot \pBel(x_{t+1})}.
\]
Allgemein gilt $\Bel(x=1)+\Bel(x=2)+\dots+\Bel(x=16)=1$ (zu jedem Zeitpunkt).
Wir haben bereits $\pBel_1(x=1)=\pBel_1(x=2)=\dots=\pBel_1(x=16)=\tfrac{1}{16}$.

\textbf{1.\ Der Roboter detektiert eine Landmarke.}

\[
  \Bel_1(x) = \eta^{-1}\cdot 0.7\cdot\tfrac{1}{16} = \eta^{-1}\cdot\tfrac{14}{320}
  \quad\text{für } x\in\{1,4,6,9,13\},
\]
\[
  \Bel_1(x) = \eta^{-1}\cdot 0.25\cdot\tfrac{1}{16} = \eta^{-1}\cdot\tfrac{5}{320}
  \quad\text{für } x\in\{2,3,5,7,8,10,11,12,14,15,16\},
\]
mit $\eta^{-1} = \dfrac{320}{5\cdot 14 + 11\cdot 5} = \dfrac{320}{125}$.

Dies bedeutet, dass nach dem ersten Schritt folgende Wsk.\ herrschen:
\[
  \Bel_1(x) =
  \begin{cases}
    \tfrac{14}{125}=0.112 & x\in\{1,4,6,9,13\},\\[0.3em]
    \tfrac{5}{125}=0.04   & x\in\{2,3,5,7,8,10,11,12,14,15,16\}.
  \end{cases}
\]

\textbf{2.\ Der Roboter bewegt sich $2$ Zellen im Uhrzeigersinn.}

Hier werden lediglich die Aufenthaltswahrscheinlichkeiten verschoben:
\[
  \pBel_2(x) =
  \begin{cases}
    0.112 & x\in\{3,6,8,11,15\},\\
    0.04  & x\in\{1,2,4,5,7,9,10,12,13,14,16\}.
  \end{cases}
\]

\textbf{3.\ Der Roboter detektiert wieder eine Landmarke.}

Hier müssen wir nun expliziter einzelne Fälle unterscheiden, um die Wahrscheinlichkeit
auszurechnen.
\[
  \Bel_2(x) = \eta^{-1}\cdot 0.7\cdot 0.04 = \eta^{-1}\cdot\tfrac{280}{10000}
  \quad\text{für } x\in\{1,4,9,13\},
\]
\[
  \Bel_2(x) = \eta^{-1}\cdot 0.25\cdot 0.04 = \eta^{-1}\cdot\tfrac{100}{10000}
  \quad\text{für } x\in\{2,5,7,10,12,14,16\},
\]
\[
  \Bel_2(x) = \eta^{-1}\cdot 0.25\cdot 0.112 = \eta^{-1}\cdot\tfrac{280}{10000}
  \quad\text{für } x\in\{3,8,11,15\},
\]
\[
  \Bel_2(x=6) = \eta^{-1}\cdot 0.7\cdot 0.112 = \eta^{-1}\cdot\tfrac{784}{10000},
\]
mit $\eta^{-1} = \dfrac{10000}{8\cdot 280 + 7\cdot 100 + 784} = \dfrac{10000}{3724}$.

Durch Einsetzen von $\eta^{-1}$ erhalten wir dadurch folgende Werte:
\[
  \Bel_2(x) =
  \begin{cases}
    \tfrac{280}{3724}\approx 0.0752 & x\in\{1,3,4,8,9,11,13,15\},\\[0.3em]
    \tfrac{100}{3724}\approx 0.0269 & x\in\{2,5,7,10,12,14,16\},\\[0.3em]
    \tfrac{784}{3724}\approx 0.211  & x=6.
  \end{cases}
\]

\textbf{4.\ Der Roboter bewegt sich $4$ Zellen im Uhrzeigersinn.}

Wie in Schritt~zwei werden hier die Wsk.\ nur verschoben:
\[
  \pBel_3(x) =
  \begin{cases}
    \tfrac{280}{3724} & x\in\{1,3,5,7,8,12,13,15\},\\[0.3em]
    \tfrac{100}{3724} & x\in\{2,4,6,9,11,14,16\},\\[0.3em]
    \tfrac{784}{3724} & x=10.
  \end{cases}
\]

\textbf{5.\ Der Roboter detektiert keine Landmarke.}

Dies verläuft ganz analog zu Schritt~drei, lediglich nun mit den Gegenwahrscheinlichkeiten:
\[
  \Bel_3(x) = \eta^{-1}\cdot 0.3\cdot\tfrac{280}{3724} = \eta^{-1}\cdot\tfrac{84}{3724}
  \quad\text{für } x\in\{1,13\},
\]
\[
  \Bel_3(x) = \eta^{-1}\cdot 0.75\cdot\tfrac{100}{3724} = \eta^{-1}\cdot\tfrac{75}{3724}
  \quad\text{für } x\in\{2,11,14,16\},
\]
\[
  \Bel_3(x) = \eta^{-1}\cdot 0.75\cdot\tfrac{280}{3724} = \eta^{-1}\cdot\tfrac{210}{3724}
  \quad\text{für } x\in\{3,5,7,8,12,15\},
\]
\[
  \Bel_3(x) = \eta^{-1}\cdot 0.3\cdot\tfrac{100}{3724} = \eta^{-1}\cdot\tfrac{30}{3724}
  \quad\text{für } x\in\{4,6,9\},
\]
\[
  \Bel_3(x=10) = \eta^{-1}\cdot 0.75\cdot\tfrac{784}{3724} = \eta^{-1}\cdot\tfrac{588}{3724},
\]
mit $\eta^{-1} = \dfrac{3724}{2\cdot 84 + 4\cdot 75 + 6\cdot 210 + 3\cdot 30 + 588} = \dfrac{3724}{2406}$.

Ergo ergibt sich durch die entsprechende Normierung:
\[
  \Bel_3(x) =
  \begin{cases}
    \tfrac{84}{2406}\approx 0.035  & x\in\{1,13\},\\[0.3em]
    \tfrac{75}{2406}\approx 0.031  & x\in\{2,11,14,16\},\\[0.3em]
    \tfrac{210}{2406}\approx 0.087 & x\in\{3,5,7,8,12,15\},\\[0.3em]
    \tfrac{30}{2406}\approx 0.012  & x\in\{4,6,9\},\\[0.3em]
    \tfrac{588}{2406}\approx 0.244 & x=10.
  \end{cases}
\]

Dies bedeutet, dass mit der höchsten Wsk.\ schlussendlich Zelle~10 erreicht wird
($\Bel_3(x=10)\approx 0.244$). Rechnet man dies mittels der Aktionen zurück, korrespondiert dies zu
Zelle~4 als wahrscheinlichster Startpunkt.

\section*{Berechnung der Aufenthaltswahrscheinlichkeit (mit unbekannter Richtung)}

Da keine Angabe über die Richtung vorliegt, gehen wir von einer $50/50$-Verteilung für die
Richtung aus. Die Prädiktion wird damit zu
\[
  \pBel(x_{t+1}) = \sum_{x_t} P(x_{t+1}\mid x_t,u_t)\cdot \Bel(x_t),
\]
wobei jede Zelle nun stets aus zwei Vorzuständen erreicht werden kann.

\textbf{1.\ Der Roboter detektiert eine Landmarke.}

Hier lassen sich die Ergebnisse einfach $1{:}1$ übertragen:
\[
  \Bel_1(x) =
  \begin{cases}
    \tfrac{14}{125}=0.112 & x\in\{1,4,6,9,13\},\\[0.3em]
    \tfrac{5}{125}=0.04   & x\in\{2,3,5,7,8,10,11,12,14,15,16\}.
  \end{cases}
\]

\textbf{2.\ Der Roboter bewegt sich $2$ Zellen.}

Grundsätzlich das gleiche Vorgehen, jedoch kann jede Zelle stets aus zwei Vorzuständen erreicht
werden (z.\,B.\ Zelle~1 und Zelle~5):
\[
  \pBel_2(x) = 0.5\cdot\tfrac{5}{125} + 0.5\cdot\tfrac{5}{125} = \tfrac{10}{250} = 0.04
  \quad\text{für } x\in\{1,5,9,10,12,13,14,16\},
\]
\[
  \pBel_2(x) = 0.5\cdot\tfrac{14}{125} + 0.5\cdot\tfrac{5}{125} = \tfrac{19}{250} = 0.076
  \quad\text{für } x\in\{2,3,4,6,7,8\},
\]
\[
  \pBel_2(x) = 0.5\cdot\tfrac{14}{125} + 0.5\cdot\tfrac{14}{125} = \tfrac{28}{250}
  \quad\text{für } x\in\{11,15\}.
\]

\textbf{3.\ Der Roboter detektiert wieder eine Landmarke.}

Hier gelten erneut die Formeln aus der vorherigen Aufgabe, also
\[
  \Bel(x_{t+1}) = \eta^{-1}\cdot P(z_t\mid x_{t+1})\cdot \pBel(x_{t+1})
\]
sowie
\[
  \eta^{-1} = \dfrac{1}{\sum_{x_{t+1}} P(z_t\mid x_{t+1})\cdot \pBel(x_{t+1})}.
\]
\[
  \Bel_2(x) = \eta^{-1}\cdot 0.7\cdot\tfrac{10}{250} = \eta^{-1}\cdot\tfrac{140}{5000}
  \quad\text{für } x\in\{1,9,13\},
\]
\[
  \Bel_2(x) = \eta^{-1}\cdot 0.25\cdot\tfrac{19}{250} = \eta^{-1}\cdot\tfrac{95}{5000}
  \quad\text{für } x\in\{2,3,7,8\},
\]
\[
  \Bel_2(x) = \eta^{-1}\cdot 0.7\cdot\tfrac{19}{250} = \eta^{-1}\cdot\tfrac{266}{5000}
  \quad\text{für } x\in\{4,6\},
\]
\[
  \Bel_2(x) = \eta^{-1}\cdot 0.25\cdot\tfrac{10}{250} = \eta^{-1}\cdot\tfrac{50}{5000}
  \quad\text{für } x\in\{5,10,12,14,16\},
\]
\[
  \Bel_2(x) = \eta^{-1}\cdot 0.25\cdot\tfrac{28}{250} = \eta^{-1}\cdot\tfrac{140}{5000}
  \quad\text{für } x\in\{11,15\},
\]
mit $\eta^{-1} = \dfrac{5000}{3\cdot 140 + 4\cdot 95 + 2\cdot 266 + 5\cdot 50 + 2\cdot 140}
= \dfrac{5000}{1862}$.

Einsetzen zum Ausrechnen von $\Bel(x_{t+1})$:
\[
  \Bel_2(x) =
  \begin{cases}
    \tfrac{140}{1862}\approx 0.0752 & x\in\{1,9,11,13,15\},\\[0.3em]
    \tfrac{95}{1862}\approx 0.051   & x\in\{2,3,7,8\},\\[0.3em]
    \tfrac{266}{1862}\approx 0.143  & x\in\{4,6\},\\[0.3em]
    \tfrac{50}{1862}\approx 0.027   & x\in\{5,10,12,14,16\}.
  \end{cases}
\]

\textbf{4.\ Der Roboter bewegt sich $4$ Zellen.}

Ähnlich zu Schritt~zwei gibt es je $2$ Zustände, aus denen man ein Feld erreichen kann
(z.\,B.\ Feld~5 aus~1 und~9):
\[
  \pBel_3(x) = 0.5\cdot\tfrac{140}{1862} + 0.5\cdot\tfrac{50}{1862} = \tfrac{190}{3724}
  \quad\text{für } x\in\{1,9\},
\]
\[
  \pBel_3(x) = 0.5\cdot\tfrac{266}{1862} + 0.5\cdot\tfrac{50}{1862} = \tfrac{316}{3724}
  \quad\text{für } x\in\{2,8,10,16\},
\]
\[
  \pBel_3(x) = 0.5\cdot\tfrac{95}{1862} + 0.5\cdot\tfrac{140}{1862} = \tfrac{235}{3724}
  \quad\text{für } x\in\{3,7,11,15\},
\]
\[
  \pBel_3(x) = 0.5\cdot\tfrac{95}{1862} + 0.5\cdot\tfrac{50}{1862} = \tfrac{145}{3724}
  \quad\text{für } x\in\{4,6,12,14\},
\]
\[
  \pBel_3(x) = 0.5\cdot\tfrac{140}{1862} + 0.5\cdot\tfrac{140}{1862} = \tfrac{280}{3724}
  \quad\text{für } x\in\{5,13\}.
\]

\textbf{5.\ Der Roboter detektiert keine Landmarke.}

Genauso wie Schritt~3, nur mit den Gegenwahrscheinlichkeiten (wie in der anderen Aufgabe):
\[
  \Bel_3(x) = \eta^{-1}\cdot 0.3\cdot\tfrac{190}{3724} = \eta^{-1}\cdot\tfrac{228}{14896}
  \quad\text{für } x\in\{1,9\},
\]
\[
  \Bel_3(x) = \eta^{-1}\cdot 0.75\cdot\tfrac{316}{3724} = \eta^{-1}\cdot\tfrac{948}{14896}
  \quad\text{für } x\in\{2,8,10,16\},
\]
\[
  \Bel_3(x) = \eta^{-1}\cdot 0.75\cdot\tfrac{235}{3724} = \eta^{-1}\cdot\tfrac{705}{14896}
  \quad\text{für } x\in\{3,7,11,15\},
\]
\[
  \Bel_3(x) = \eta^{-1}\cdot 0.3\cdot\tfrac{145}{3724} = \eta^{-1}\cdot\tfrac{174}{14896}
  \quad\text{für } x\in\{4,6\},
\]
\[
  \Bel_3(x) = \eta^{-1}\cdot 0.75\cdot\tfrac{145}{3724} = \eta^{-1}\cdot\tfrac{435}{14896}
  \quad\text{für } x\in\{12,14\},
\]
\[
  \Bel_3(x=5) = \eta^{-1}\cdot 0.75\cdot\tfrac{280}{3724} = \eta^{-1}\cdot\tfrac{840}{14896},
  \qquad
  \Bel_3(x=13) = \eta^{-1}\cdot 0.3\cdot\tfrac{280}{3724} = \eta^{-1}\cdot\tfrac{336}{14896},
\]
mit $\eta^{-1} = \dfrac{14896}{2\cdot 228 + 4\cdot 948 + 4\cdot 705 + 2\cdot 174 + 2\cdot 435 + 840 + 336}
= \dfrac{14896}{9462}$.

Demnach sind schlussendlich die Wsk.:
\[
\begin{aligned}
  &P(1\mid\bar J)=P(9\mid\bar J)\approx 0.024, \quad
  P(2\mid\bar J)=P(10\mid\bar J)=P(8\mid\bar J)=P(16\mid\bar J)\approx 0.1, \\[0.2em]
  &P(3\mid\bar J)=P(11\mid\bar J)=P(7\mid\bar J)=P(15\mid\bar J)\approx 0.074, \quad
  P(4\mid\bar J)=P(6\mid\bar J)\approx 0.018, \\[0.2em]
  &P(12\mid\bar J)=P(14\mid\bar J)\approx 0.046, \quad
  P(5\mid\bar J)\approx 0.089, \quad
  P(13\mid\bar J)\approx 0.035.
\end{aligned}
\]
Bzw.\ mit den entsprechenden Brüchen:
\[
  \Bel_3(x) =
  \begin{cases}
    \tfrac{228}{9462}\approx 0.024 & x\in\{1,9\},\\[0.3em]
    \tfrac{948}{9462}\approx 0.1   & x\in\{2,8,10,16\},\\[0.3em]
    \tfrac{705}{9462}\approx 0.075 & x\in\{3,7,11,15\},\\[0.3em]
    \tfrac{174}{9462}\approx 0.018 & x\in\{4,6\},\\[0.3em]
    \tfrac{435}{9462}\approx 0.046 & x\in\{12,14\},\\[0.3em]
    \tfrac{840}{9462}\approx 0.089 & x=5,\\[0.3em]
    \tfrac{336}{9462}\approx 0.036 & x=13.
  \end{cases}
\]

Hier sind am wahrscheinlichsten die Zellen $2$, $8$, $10$ oder $16$ als Endposition
(mit $\Bel_3(x=2)=\Bel_3(x=8)=\Bel_3(x=10)=\Bel_3(x=16)\approx 0.1$). Jedoch ist alles allgemein
sehr unklar geworden, nicht so eindeutig wie bei der anderen Aufgabe. Eine Rückrechnung auf den
„wahrscheinlichsten“ Initialzustand wäre zwar grundsätzlich möglich, würde allerdings aufgrund der
gewichteten Kanten mit der zusätzlichen Komplexität den Rahmen etwas sprengen.

\end{document}
